\section{Changed conversion terms change the accounting question}
\label{sec:financing-changes}

A maturity conversion, an early repurchase and an inducement to convert all
remove or alter debt, but they do not allocate the resulting consideration in
the same way. The issuer must retain the old conversion rights, the revised
rights, the date of amendment and the actual settlement. A record that says
only \emph{converted} loses the facts needed to distinguish these events.
The following cases use the retained Australian standards for explicit
calendar-2024/2025 for-profit Tier~1 illustrations, without optional later
amendments. They are pretax: tax consequences remain uncomputed.

For an early repurchase with unchanged conversion privileges, the issuer
allocates the consideration between the liability and equity components using
the method of the original separation. The liability allocation uses its
current fair value; the difference from its carrying amount enters profit or
loss, while the equity allocation stays in equity
\citep[AG33--34 and IE39--46]{aasb132}. The standard's example has a
face amount of 1,000, half-yearly coupons of 50 and ten payments remaining.
The original half-year discount rate is 5.5\%; the contemporaneous comparable
nonconvertible debt rate is 4\%. For either rate $r$, the remaining debt value is
\begin{equation}
 V(r)=50\frac{1-(1+r)^{-10}}{r}+1{,}000(1+r)^{-10}.
 \label{eq:changed-financing-pv}
\end{equation}
The first term values the ten coupons; the second values principal at the
last payment. The source rounds the carrying and current fair values to
962 and 1,081. Repurchase consideration of 1,700 then allocates 619 to equity
and produces debt settlement expense of 119. The entries debit debt 962,
expense 119 and equity 619, and credit cash 1,700. Expensing the whole
$1{,}700-962=738$ would misclassify the equity purchase. The original equity
component of 60 remains within equity; it does not cap the 619 deduction.
Exact present-value fractions and the source's rounded presentation remain
separate in the accompanying calculations.

An inducement asks a different valuation question. At the amendment date,
compare the fair value of what the holder receives under the revised
conversion terms with what the holder would have received under the original
terms. Their difference is an inducement loss
\citep[AG35 and IE47--50]{aasb132}. In the retained example, changing the
conversion price from 25 to 20 on face amount 1,000 increases shares from
40 to 50. At the amendment-date share quotation of 40, the additional ten
shares are worth 400. A later quotation of 50 would give 500 and would use
the wrong measurement date.

For an explicitly immediate, unconditional share-only conversion, suppose the
liability's admitted carrying amount is 960 and the original equity component
is 60. This carrying amount is a separate illustrative input. The resulting
entry debits debt 960 and inducement expense 400, and credits issued equity
1,360; 50 shares are issued. Original conversion equity remains in equity.
The 400 expense offsets the corresponding part of the issued equity in total
equity, while derecognizing the debt increases net equity by 960. There is no
cash movement. Noncash financing disclosure preserves the event without
inventing a financing receipt \citep[43--44]{aasb107dated}.

This immediate case does not determine the recognition timing of a delayed,
contingent or cash inducement. Those terms must first establish the obligation.
Nor does an ability to settle in shares establish equity classification.
A derivative with a net-cash alternative fails the all-equity-alternatives
condition, and an obligation to deliver a variable number of shares worth a
fixed sum does not represent a residual ownership interest
\citep[21 and 26--27]{aasb132}. The system must refuse the fixed-for-fixed
method when these facts are missing or contradicted, while retaining the
instrument and its unresolved classification.

\section{A hedge reserve is neither cash nor a second gain}
\label{sec:hedge-lifecycle}

Suppose a company hedges a highly probable inventory purchase. A derivative
gain can arrive before the inventory affects earnings. Cash-flow hedge
accounting preserves that timing through a separate equity reserve, but only
after the relationship qualifies. The designation must identify an eligible
instrument and item, the hedged risk, the risk-management objective and the
assessment method at inception. An economic relationship, the absence of
credit-risk dominance and an appropriate actual hedge ratio are separate
requirements \citep[6.2--6.4 and B6.4.1--19]{aasb9dated}. Favorable price
correlation alone does not establish them. Unknown designation or valuation
information remains unresolved; it cannot be repaired by assigning zero or
placing a gain in other comprehensive income.

Take one documented, offsetting exposure and a derivative worth zero at
designation. The illustration includes the whole derivative, no excluded
option or forward component and no intermediate reserve release. Let $D_t$
be its cumulative gain, including settled gains, and $H_t$ the cumulative
change in the hedged item's present value attributable to the designated
risk. Their signs refer to their economic values. With admitted offsetting
changes, the cumulative effective gain $E_t$ and ineffectiveness $I_t$ are
\begin{equation}
 E_t=\operatorname{sgn}(D_t)\min\{|D_t|,|H_t|\},
 \qquad I_t=D_t-E_t.
 \label{eq:hedge-cumulative-lower}
\end{equation}
The reserve follows the lower absolute cumulative amount; effective changes
enter other comprehensive income and the remainder enters profit or loss
\citep[6.5.11]{aasb9dated}. This formula handles the stated offsetting case;
a complex exposure outside it needs its own measurement method. In particular,
one same-direction period is not itself proof that no economic relationship
exists. Item valuation must include time value and only the item's relevant
features; a hypothetical derivative cannot import features that exist solely
in the actual hedging instrument \citep[B6.4.5 and B6.5.4--6]{aasb9dated}.

Initially, supplied values $D_1=12$ and $H_1=-10$ produce effective gain 10
and profit 2. At the next measurement, $D_2=15$ and $H_2=-14$ produce
effective gain 14 and cumulative profit 1. The second period therefore records
a derivative increase of 3, other comprehensive income of 4 and a loss of 1.
Applying the minimum only to that period's changes would incorrectly give
effective gain 3. The retained cumulative values are essential economic memory.

Now buy the inventory for 114 and settle the derivative for cash of 15.
Settlement removes the derivative asset; it does not by itself release the
14 reserve. Because the forecast transaction has produced a nonfinancial
asset, debit the reserve 14 and credit inventory 14. Inventory cost becomes
100. This direct basis adjustment is not an other-comprehensive-income
reclassification \citep[6.5.11(d)(i)]{aasb9dated}. Selling the inventory for
130 subsequently produces revenue 130 and cost of sales 100.

With opening cash and equity of 200 and no other movements, the full sequence
has cumulative pretax profit 31, other comprehensive income 14 and a separate
equity-to-inventory basis transfer of $-14$. Closing cash and pretax equity
are both 231. The equity reconciliation is
\begin{equation}
 200+31+14-14=231,
 \label{eq:hedge-equity-bridge}
\end{equation}
where omitting the direct basis transfer would overstate equity. Gross
operating cash is $-114+15+130=31$: the identified hedge's cash flows follow
the inventory position's classification \citep[16]{aasb107dated}. A matching
terminal cash total does not excuse a wrong interim reserve, profit or basis.

If the hedged payments are interest, the reserve instead enters earnings when
the interest affects earnings \citep[6.5.11(d)(ii)]{aasb9dated}. With an
effective gain of 8, derivative settlement of 8 and cash interest of 18,
release 8 from the reserve into earnings. Expense is then 10. The release
has an OCI reclassification of $-8$, but no second cash receipt. The
illustrative nonfinancial company consistently elects operating classification
for interest under the retained rule; net operating cash is $-10$
\citep[31--33]{aasb107dated}. A later presentation regime needs its own
applicability decision.

Stopping hedge accounting also requires a reason. If the future cash flows
remain expected, the accumulated reserve remains until the flows occur,
subject to the unrecoverable-loss rule. If they are no longer expected, the
reserve is reclassified immediately to profit or loss. No longer highly
probable is not the same as no longer expected
\citep[6.5.11(d)(iii) and 6.5.12]{aasb9dated}. The worked alternatives keep
a gain reserve of 10 in the first case and release it in the second. A loss
reserve of 10 with an unrecoverable portion of 4 instead releases expense 4
and leaves a loss reserve of 6. Each decision changes a named exposure and
reserve; it cannot be inferred from the derivative's cash settlement alone.

\subsection{Fair-value adjustments and benchmark changes}

An admitted fair-value hedge of amortised-cost debt uses a different route.
Both the derivative change and the risk-attributable debt adjustment enter
profit or loss; the latter changes the debt carrying amount
\citep[6.5.8]{aasb9dated}. For debt initially carried at 100, a derivative
gain of 8 and debt-value increase of 10 give a derivative asset of 8, debt
of 110 and net loss of 2. No cash-flow hedge reserve is created.

Once the debt ceases receiving hedge adjustments, the retained adjustment
must be amortized using a recalculated effective interest rate
\citep[6.5.10]{aasb9dated}. Suppose contractual principal of 100 and coupon
of 21 are payable together one year later. From carrying value 110, the
recalculated rate is $121/110-1=10\%$, giving effective interest expense
11 before settlement. Cash interest remains 21 and principal remains 100;
the difference between cash interest and effective expense amortizes the
10 adjustment. Reusing the old 21\% rate would over-accrete the liability.
Missing contractual dates or payments leave that calculation unresolved.

Benchmark reform is not a general exemption from this analysis. The retained
standard's rate-update expedient applies to necessary reform changes whose
new cash-flow basis is economically equivalent to the old basis. Additional
changes then receive the ordinary modification or derecognition analysis
\citep[5.4.5--9]{aasb9dated}. Comparing one 4\% coupon with 3.8\% plus
a 0.2\% spread verifies one number, not equivalence of all future terms.
The amendment record must retain the legal change, rates, spreads, reset
terms, dates and the supporting equivalence analysis.

A qualifying change to formal hedge designation is neither automatic
discontinuation nor a new relationship. Existing cash-flow reserves follow
the alternative benchmark under the specified conditions rather than being
reset to zero. Other qualifying requirements continue to apply; additional
changes and mixed benchmark groups need their own assessment
\citep[6.9.1--10]{aasb9dated}. The rate-specific 24-month exception for a
separately identifiable risk component likewise requires a supported
expectation and ends on the stated adverse reassessment
\citep[6.9.11--13]{aasb9dated}.

The exact conditional journals and adverse alternatives are retained in
\path{data/accounting-master/finite-method-decisions-20261015/financing/};
the field and event specification is
\path{docs/design/proposal/methods/financing-changes-20261015.md}.
These cases specify the financial consequences of admitted values and
eligibility findings. The next section extends the treatment to excluded
components and rebalancing. Actual valuation, delayed or contingent inducements
and group hedges retain their named method and evidence obligations.
Supplying a market value is useful input; it is not a completed valuation method.
